\documentclass{article}
\usepackage[T1]{fontenc}
\usepackage{lmodern}
\usepackage{graphicx}
\usepackage{amsmath,amssymb}
\usepackage[a4paper,margin=2cm]{geometry}
\usepackage[absolute,overlay]{textpos}
\setlength{\TPHorizModule}{1mm}
\setlength{\TPVertModule}{1mm}
\usepackage[indonesian]{babel}
\usepackage{hyperref}
\setlength{\emergencystretch}{2em}
% unit: Halaman 35

\begin{document}

% === Halaman 35 (llm) ===

\section*{Mode DES}

\begin{itemize}
    \item DES dapat dioperasikan dengan mode ECB, CBC, OFB, CFB, dan mode \textit{counter}.
    \item Namun karena kesederhanaannya, mode ECB lebih sering digunakan pada paket komersil, sehingga dikenal dengan nama DES-ECB.
\end{itemize}

\end{document}
